\begin{abstract}
We characterize minimax recovery of the treated causal response under a global propensity-tail restriction. The model combines consistency and conditional exchangeability with a covariate density bounded below on a cube, intrinsic Hölder response smoothness of known order \(\beta>1\), a bounded treated conditional mean, and uniformly conditional sub-Gaussian residuals. If \(d\) is the covariate dimension, \(n\) is the sample size, and \(\gamma>1\) is the overlap exponent, the pointwise and expected spatial-supremum minimax risks have scale
\[
n^{-\beta/(2\beta+d\gamma/(\gamma-1))}.
\]
An explicit estimator fits local polynomials with equal aggregate weight across fixed template microcells and selects its mesh through treated observation counts. With model constants fixed, its expected spatial-supremum guarantee attains this pure-power rate uniformly over fixed compact overlap intervals with distinct endpoints in \((1,\infty)\), using the known smoothness and observed treatment design. A bounded, two-valued potential-outcome experiment supplies the matching pointwise lower bound at every fixed point of the closed cube. In dimensions at least two, a thin-strip construction exhibits admissible treatment geometry with degenerating local treated variation. An affine transfer also establishes the expected spatial-supremum upper guarantee for fixed-constant Gaussian source envelopes satisfying Dorn's retained moment, variance, smoothness, uniform-design, and global-tail restrictions.
\end{abstract}

\section{Introduction}

This paper establishes the minimax rate for recovering a treated causal response throughout a covariate cube when treatment availability is governed by a global propensity tail. Under consistency and conditional exchangeability, the target is the conditional mean of the treated potential outcome, identified through treated outcome regression. We study laws with a covariate density bounded below, a response in an intrinsic cube Hölder ball of known order, a uniformly bounded treated conditional mean, and uniformly conditional sub-Gaussian residuals. Write \(n\) for sample size, \(\beta>1\) for response smoothness, \(d\ge1\) for covariate dimension, and \(\gamma>1\) for the overlap exponent. Within this fixed-constant class, one explicit estimator attains the expected spatial-supremum rate
\[
n^{-\beta/(2\beta+d\gamma/(\gamma-1))}.
\]
The upper guarantee is uniform over any fixed compact exponent interval with distinct endpoints contained in \((1,\infty)\). For each fixed exponent and each fixed evaluation point, a matching lower bound establishes the same minimax scale for pointwise expected absolute error.

Response recovery is a basic component of causal adjustment in the potential-outcome framework \citep{rubin1974,rosenbaum1983,imbens2015}. Its statistical difficulty depends on where treated observations are available. Our overlap restriction bounds the covariate probability of low-propensity regions by a power of the propensity threshold. This restriction quantifies the prevalence of treatment scarcity while accommodating measurable propensities and heterogeneous local treatment geometry. The target remains the continuous treated response over the entire closed cube, and the spatial loss measures the largest estimation error within each sample realization before taking expectation.

The construction separates polynomial conditioning from treatment-count precision. Each partition cube contains scaled copies of a fixed collection of small microcells around polynomial interpolation nodes. The fit gives every microcell the same aggregate weight, averaging treated observations within each cell. This weighting keeps the empirical Gram matrix close to a fixed positive-definite template whenever the cells are occupied. A count selector chooses the finest candidate mesh at which every microcell contains enough treated observations relative to the approximation scale. The resulting estimator uses the observed sample, sample size, dimension, known smoothness, and response bound, as specified in \cref{obj:def:mesh-selector,obj:def:estimator}.

The statistical explanation rests on an ordered information profile. The density lower bound assigns covariate mass to each template microcell, and the global propensity tail controls the treated mass of unions of disjoint microcells. Together, they imply a polynomial lower bound in rank for the weakest treated-cell masses across partition cubes. This profile distinguishes the least-informed cubes from the progressively better-informed cubes. Combining its rank-dependent stochastic control with the common scale cap supplied by count feasibility yields the pure-power expected spatial-supremum guarantee in \cref{obj:thm:adaptive-minimax}. Local polynomial approximation and deterministic template conditioning complete the upper bound.

The converse uses a bounded causal experiment. For every fixed evaluation point, including points on faces and corners, a pair of laws places a smooth response perturbation near that point and shares a radial propensity satisfying the global tail envelope. The laws admit potential outcomes taking values in a common two-point set. Their response separation and bounded relative entropy give the pointwise lower bound in \cref{obj:lem:bounded-lower-pair}. Thus the minimax characterization applies to the full conditional sub-Gaussian response class, with the converse already supported by a bounded subclass.

The closest theorem-level comparison is \citet[Assumptions 1--3 and Theorem 1(ii), p.~12]{Dorn2026GlobalOverlap}. For a fixed Gaussian uniform-design setup satisfying Dorn's retained restrictions and local anti-concentration condition, Dorn establishes a probability upper bound at the same pure-power scale in an essential-supremum criterion. The source describes a count-based estimator computable without the overlap exponent, while its probability upper guarantee fixes the source family, its common constants, and that exponent before selecting the estimator and threshold constant \citep[Lemmas 11 and 13, pp.~23--24; proof, pp.~27--28]{Dorn2026GlobalOverlap}. That construction provides a direct antecedent for mesh selection. Dorn's local anti-concentration condition additionally links every sufficiently small neighborhood's largest selected propensity to a common positive fraction of its covariate mass, with constants common across the fixed family and its selected propensity representatives. Our expected spatial-supremum guarantee is established over the global-tail class in \cref{obj:def:model}, using one estimator and common constants across the stated compact exponent range.

The geometric and source-model comparisons make this distinction precise. In dimensions at least two, \cref{obj:prop:strict-enlargement} constructs a thin-strip law within the causal class whose displayed propensity violates the selected-version local anti-concentration condition and whose normalized treated-design second-coordinate moment tends to zero. Separately, \cref{obj:prop:dorn-a3-free-corollary} transfers the upper guarantee to fixed-constant Gaussian envelopes satisfying the retained source restrictions from \citet[Section 2, Assumptions 1--2, pp.~3--7]{Dorn2026GlobalOverlap}. Affine transport and fixed-cube smoothness completion supply the common full Hölder radius. The transferred guarantee controls expected spatial-supremum error uniformly over those envelopes and the compact exponent range. Its Gaussian geometric witnesses use constants selected separately for each exponent.

Related regression results characterize other aspects of heterogeneous information. \citet[Assumption M and Theorem 1, pp.~3--4]{Gaiffas2005DegenerateRates} derives pointwise rates under locally symmetric, regularly varying design densities in a one-dimensional Gaussian experiment. \citet[Assumption D, equation (2.1), and Theorems 1--2]{gaiffas2009} and \citet[Sections 1 and 3]{schmidthieber2024} retain spatial variation through location-dependent loss normalization. Our criterion instead evaluates expected unweighted spatial-supremum error over a global overlap envelope. The causal target also differs from the smooth treatment contrast studied under strict overlap by \citet[Theorems 1--2 and Remark 11]{KennedyBalakrishnanRobinsWasserman2024}: here smoothness is imposed directly on the treatment-specific response.

An appendix documents the precise scope of the mathematical results.

The exposition proceeds through related work, the model and risk criteria, and the estimator and recovery guarantee. It then develops the local treatment-geometry comparison and Gaussian transfer, followed by discussion and future work. The appendices present auxiliary identification, upper-bound, testing, geometry, and source-comparison arguments, followed by the complete main-result proofs.
